% ============================================================
% DRAFT — proposed edits to chapters/literature.tex from Theart's
% walkthrough of Chapter 2 (recorded feedback, 20 Oct 2026).
%
% This file is a draft for review, not a finished chapter section.
% Each block below is: a short header comment naming the anchor
% text in the CURRENT literature.tex, a note on what feedback
% point it answers, and the replacement text. Nothing here has
% been written into the actual chapter file — copy over whatever
% you're happy with, section by section, and edit the rest.
%
% Feedback points covered, roughly in the order Theart raised them:
%   E1. sec:grounding opens abruptly with "SayCan and Code as
%       Policies set the standard way..." — needs a grounding
%       paragraph, and shouldn't imply these are the only two
%       methods.
%   E2. SayCan's affordance/value function is used but never
%       explained: where does p(c|s,l) actually come from, how
%       is it trained?
%   E3. Code as Policies is vague about what the LLM's code
%       generation is actually conditioned on (a prompt-level API
%       + examples), both where CaP is first introduced
%       (sec:llm_planners) and the shorter repeat in sec:grounding.
%   E4. "This project's mid-mission execution follows the same
%       loop" / "this project's capability vocabulary in
%       planner.py" — these are Chapter 3 design claims sitting in
%       Chapter 2. Deferred to a forward-reference instead.
%   E5. Quote direction: two straight "..." pairs need to become
%       LaTeX `` '' pairs (lines 64 and 128 in the current file).
%       I didn't find any others in literature.tex or
%       introduction.tex — grep -n '"' both files to confirm.
%   E6. KnowNo's "raw confidence score" is passed through
%       conformal prediction, but where does the score itself come
%       from? (It's not a separate trained model — see E2's answer
%       for why this matters.)
%   E7. Section~\ref{sec:local_models}'s opening overclaims: none
%       of Sections 2.1-2.3's methods *require* a cloud-hosted LLM,
%       they were just evaluated with one.
%   E8. sec:hallucination_lit conflates two distinct failure modes
%       under "reasoning degrades over a long chain of steps": a
%       context-window capacity limit vs. one-shot-vs-reasoning
%       answer quality. Split into two, and renumbered the "second
%       problem" (now third) later in the same section + added a
%       short gap statement in the Research Gap section about the
%       un-addressed capacity limit.
%   E9. Awkward phrasing at the end of 2.4 (now 2.4's closing
%       paragraph) — reworded for clarity.
%   E10. STT section (2.5) trimmed, and rebalanced towards stating
%        the solution (in-domain/in-accent fine-tuning,
%        digiphyte/fluister-turbo) rather than only the problem —
%        with the hedge Theart gave from his own experience (a
%        full-weight general Whisper model sometimes beats it).
%
% Also includes one unprompted grammar catch (A0): a stray period
% breaks a sentence mid-clause in the very first SayCan paragraph.
% ============================================================


% ---------- A0. Stray period, sec:llm_planners (first SayCan paragraph) ----------
% CURRENT (around line 18-19):
%   ...but each candidate is re-scored by a separately learned affordance function.
%   that estimates how likely that skill is to succeed given the robot's current state. The action
% REPLACE WITH:

...but each candidate is re-scored by a separately learned affordance function
that estimates how likely that skill is to succeed given the robot's current state. The action


% ---------- E1. sec:grounding opening — needs a grounding paragraph + hedge ----------
% CURRENT (lines 56-60):
%   \mysection{Grounding plans in a bounded action set}{Grounding Plans in a Bounded Action Set}
%   \label{sec:grounding}
%
%   SayCan and Code as Policies set the standard way to connect an LLM's plan to what a robot can
%   actually do. SayCan~\cite{ahn2022saycan} pairs the LLM with a separately learned value function for
% REPLACE WITH:

\mysection{Grounding plans in a bounded action set}{Grounding Plans in a Bounded Action Set}
\label{sec:grounding}

Section~\ref{sec:llm_planners} introduced SayCan, Code as Policies and Inner Monologue at a high
level, as three systems that connect an LLM's plan to what a robot can actually do. This section
looks more closely at how the first two of those three perform that grounding step mechanically,
since the way each one scores or constrains a candidate action is what the rest of this chapter
builds on. At least two prominent methods set a standard way of doing this: SayCan~\cite{ahn2022saycan}
pairs the LLM with a separately learned value function for


% ---------- E2. How is SayCan's affordance value function obtained? ----------
% INSERT as a new paragraph right after the SayCan equation block, i.e. after
% the current line ending "...and the 0.6/0.2/0.1/0.8 values are
% $p(c_\pi \mid s, \ell_\pi)$." and before "Code as Policies~\cite{liang2023codeaspolicies}
% solves the same problem a different way: ...".

Ahn et al.\ do not get $p(c_\pi \mid s, \ell_\pi)$ for free. Each skill's affordance value function
is trained separately from the language model, as a reinforcement-learning critic fitted on
logged robot trajectories: many attempts at that skill are recorded from a range of starting
states, each ending in a success or failure outcome, and the critic is trained to predict that
outcome's probability from the observation alone. Obtaining $p(c_\pi \mid s, \ell_\pi)$ is
therefore not a property of the language model at all; it requires a separate, per-skill dataset
of practiced attempts and their outcomes before any value function exists to score against.


% ---------- E3a. Code as Policies vagueness, sec:llm_planners (first mention) ----------
% CURRENT (lines 26-32):
%   Code as Policies \cite{liang2023codeaspolicies} takes a different route. Rather than picking
%   from a fixed skill library, the LLM writes code, calling perception functions, control
%   primitives, and other functions it has been shown examples of, composed however the
%   instruction requires. This lets it handle instructions no one anticipated in advance and
%   assign precise parameters (a velocity, a threshold) to descriptions that were only ever given
%   in words, at the cost of needing a language model capable of writing correct, composable code
%   rather than just picking from a short list.
% REPLACE WITH:

Code as Policies \cite{liang2023codeaspolicies} takes a different route. Rather than picking
from a fixed skill library, the LLM writes code against a small, hand-written API of perception
functions, control primitives and helper utilities, whose signatures and example usage are given
to the model directly in its prompt, alongside a handful of worked examples pairing an instruction
with the code that solves it. Composing calls to only that API, in whatever order and nesting the
instruction requires, is what lets it handle instructions no one anticipated in advance and assign
precise parameters (a velocity, a threshold) to descriptions that were only ever given in words,
at the cost of needing a language model capable of writing correct, composable code against that
API rather than just picking from a short list.


% ---------- E3b. Code as Policies, sec:grounding (shorter repeat mention) ----------
% CURRENT (around what is currently lines 106-108):
%   Code as Policies~\cite{liang2023codeaspolicies} solves the same problem a different way: the LLM
%   writes Python that calls a fixed set of perception and control functions directly. If the code does
%   not compile against that API, the step cannot run.
% REPLACE WITH:

Code as Policies~\cite{liang2023codeaspolicies} solves the same problem a different way: as
Section~\ref{sec:llm_planners} described, the LLM writes Python against the fixed API of
perception and control functions given to it in its prompt. If the code does not compile against
that API, the step cannot run.


% ---------- E4a. "capability vocabulary in planner.py" — Ch3 design claim in Ch2 ----------
% CURRENT (around what is currently lines 110-112):
%   Both approaches give the planner a fixed set of actions it cannot go outside of. This project's
%   capability vocabulary in \texttt{planner.py} plays the same role: a step whose target or action
%   falls outside that vocabulary is rejected rather than passed through to the rover.
% REPLACE WITH:

Both approaches give the planner a fixed set of actions it cannot go outside of, a property
Chapter~\ref{chap:design} returns to when this project's own capability vocabulary is introduced.


% ---------- E4b. "mid-mission execution follows the same loop" — Ch3 design claim in Ch2 ----------
% CURRENT (around what is currently lines 114-121):
%   Neither approach checks what happens after a step runs. Inner Monologue~\cite{huang2023innermonologue}
%   adds that check on top of a SayCan-style setup: after each action, the robot gets an observation
%   back, a success or failure signal, a scene description, or something the human said, and the next
%   step is planned using that observation instead of the original plan alone. Plan a step, act, observe,
%   replan. This project's mid-mission execution follows the same loop: each frame sent over
%   \texttt{/ws/execution} is what the plan gets checked against after each step, not only before the
%   mission starts, allowing the plan to respond to drift the rover has actually accumulated rather than
%   committing to a path decided before departure.
% REPLACE WITH:

Neither approach checks what happens after a step runs. Inner Monologue~\cite{huang2023innermonologue}
adds that check on top of a SayCan-style setup: after each action, the robot gets an observation
back, a success or failure signal, a scene description, or something the human said, and the next
step is planned using that observation instead of the original plan alone. Plan a step, act,
observe, replan. Chapter~\ref{chap:design} returns to this plan-act-observe-replan loop when
motivating this project's own mid-mission execution design.


% ---------- E5. Quote direction fixes ----------
% Line 64, CURRENT:
%   with a concrete instruction, "how would you put an apple on the table?" The LLM alone ranks
% REPLACE WITH:

with a concrete instruction, ``how would you put an apple on the table?'' The LLM alone ranks

% Line 128, CURRENT:
%   observations rule out every other reading of it. "Bring the bowl" is unambiguous with one bowl on the counter
% REPLACE WITH:

observations rule out every other reading of it. ``Bring the bowl'' is unambiguous with one bowl on the counter


% ---------- E6. KnowNo's "raw confidence score" — where does it come from? ----------
% CURRENT (around what is currently lines 136-140):
%   \citet{ren2023knowno} give active modification a statistical guarantee, in a method called KnowNo, and \citeauthor{hori2025longhorizon} cite it directly as a working example of the approach.
%   KnowNo does not let the planner commit to its single highest-scoring next step. The planner's raw confidence scores for each candidate step are passed through a conformal prediction layer,
%   which converts them into a calibrated set of steps that are all plausible enough to keep. A singleton set means the rover proceeds on its own. A set with more than one step means the instruction was ambiguous
%   given what the rover can observe, and the rover asks a clarifying question instead of guessing.
% REPLACE WITH:

\citet{ren2023knowno} give active modification a statistical guarantee, in a method called
KnowNo, and \citeauthor{hori2025longhorizon} cite it directly as a working example of the
approach. KnowNo does not let the planner commit to its single highest-scoring next step. The
planner's raw confidence scores for each candidate step are passed through a conformal prediction
layer. These raw scores are not a separately trained model: Ren et al.\ take them straight from
the planning LLM itself, normalising its softmax probability over the candidate next steps
(optionally combined with a SayCan-style affordance score, Equation~\ref{eq:saycan_selection},
when one is available), so KnowNo adds no new model to train, only a calibration step on top of
whatever scoring mechanism the base planner already produces. Conformal prediction then converts
these scores into a calibrated set of steps that are all plausible enough to keep. A singleton set
means the rover proceeds on its own. A set with more than one step means the instruction was
ambiguous given what the rover can observe, and the rover asks a clarifying question instead of
guessing.


% ---------- E7. sec:local_models opening overclaims cloud-hosting as a requirement ----------
% CURRENT (around what is currently lines 184-189):
%   \mysection{Locally hosted models}{Locally hosted and quantised models}
%   \label{sec:local_models}
%
%   Sections~\ref{sec:llm_planners} to \ref{sec:clarification_lit} assume the rover can call out to a large, cloud-hosted LLM. A rover cannot always guarantee that access: network connectivity is not
%   reliable in every test environment, and the round-trip latency to a cloud endpoint works against the response time a voice-commanded system needs to feel usable. Running the planner locally,
%   on hardware the rover carries itself, removes both problems, but restricts the model to whatever fits in on-board memory.
% REPLACE WITH:

\mysection{Locally hosted models}{Locally hosted and quantised models}
\label{sec:local_models}

None of the grounding, coding or clarification mechanisms in Sections~\ref{sec:llm_planners} to
\ref{sec:clarification_lit} assume cloud hosting as a requirement; each needs an LLM, not
specifically a large, cloud-hosted one. In practice, however, every one of them was evaluated in
its original paper using a large model accessed over the internet, so how a smaller, locally
hosted model performs on the same task is not something that body of work establishes either way.
A rover cannot always guarantee cloud access regardless: network connectivity is not reliable in
every test environment, and the round-trip latency to a cloud endpoint works against the response
time a voice-commanded system needs to feel usable. Running the planner locally, on hardware the
rover carries itself or a machine on the same network, removes both problems, but restricts the
model to whatever fits in on-board or local memory, with an expected drop in capability relative
to the cloud-scale models these methods were evaluated with.


% ---------- E8a. Conflated failure modes, sec:hallucination_lit opening ----------
% CURRENT (around what is currently lines 218-221):
%   The methods discussed in Sections~\ref{sec:llm_planners} to \ref{sec:clarification_lit} treat each planning step largely on its own.
%   As a task plan grows longer and more compositional, two further problems appear that a single-step clarification mechanism does not cover: the LLM's reasoning can degrade over a long chain of steps,
%   producing a confidently stated plan that is simply wrong, and the rover can encounter a scene it was never prepared for, one that no single faulty component causes
%   but that the plan as a whole fails to account for.
% REPLACE WITH:

The methods discussed in Sections~\ref{sec:llm_planners} to \ref{sec:clarification_lit} treat each
planning step largely on its own. As a task plan grows longer and more compositional, three
further problems appear that a single-step clarification mechanism does not cover. The first is a
capacity limit rather than a reasoning one: a model's effective context window, how much of the
conversation and scene history its architecture and memory budget let it hold onto at once,
eventually forces it to lose track of earlier parts of a long mission, independent of how well it
reasons about what it can still see. The second is a reasoning-quality problem: scoring a
candidate action directly, in a single pass, is less reliable than first reasoning explicitly
about why that action is or is not appropriate, and skipping that reasoning step is what produces
a confidently stated plan that is simply wrong. The third is that the rover can encounter a scene
it was never prepared for, one that no single faulty component causes but that the plan as a whole
fails to account for.


% ---------- E8b. Renumber "first problem" -> "second", liang2024introspective paragraph ----------
% CURRENT, first sentence of the next paragraph:
%   \citet{liang2024introspective} address the first problem by extending the conformal prediction procedure from Section~\ref{sec:clarification_lit} with a retrieval step. Rather than scoring each candidate plan from the LLM's raw confidence alone,
% REPLACE WITH:

\citet{liang2024introspective} address the second of these, the reasoning-quality problem, by
extending the conformal prediction procedure from Section~\ref{sec:clarification_lit} with a
retrieval step. Rather than scoring each candidate plan from the LLM's raw confidence alone,


% ---------- E8c. Renumber "second problem" -> "third", AESOP/semantic-anomaly paragraph ----------
% CURRENT:
%   The second problem, an unanticipated scene rather than an unreliable plan,
%   is addressed by \citet{sinha2024aesop} and \citet{elhafsi2023semantic} from the
%    direction of runtime monitoring rather than planning.
% REPLACE WITH:

The third problem, an unanticipated scene rather than an unreliable plan,
is addressed by \citet{sinha2024aesop} and \citet{elhafsi2023semantic} from the
 direction of runtime monitoring rather than planning.


% ---------- E8d. New gap statement for the un-addressed context-window problem ----------
% INSERT in the Research Gap section (sec:gap), as a new paragraph. Current flow has
% a paragraph ending "...is not addressed by either strand of work." followed by a
% paragraph starting "Similarly, none of the planning or clarification methods...".
% Insert the new paragraph BETWEEN those two.

A related gap concerns compositional length directly. Section~\ref{sec:hallucination_lit}
distinguished a model's context-window capacity, how much mission history it can hold onto at
once, from the quality of its reasoning once given everything it needs. Introspective planning
and the anomaly-detection methods in that section address the second; none address the first. How
many steps a long-horizon command can be chained into before context-window capacity, rather than
reasoning quality, becomes the limiting factor is not quantified by any of the work reviewed in
this chapter.


% ---------- E9. Awkward phrasing, closing paragraph of sec:hallucination_lit ----------
% CURRENT:
%   Together, these three methods extend the single-step clarification of
%   Section~\ref{sec:clarification_lit} along a different axis: introspective planning
%   improves the confidence estimate feeding a single decision, while AESOP and semantic anomaly
%   detection catch failures a confidence score never gets the chance to be wrong about, because the
%    plan itself is being carried out against a scene the planner misjudged.
% REPLACE WITH:

Together, these three methods extend the single-step clarification of
Section~\ref{sec:clarification_lit} along a different axis: introspective planning sharpens the
confidence estimate feeding a single decision, while AESOP and semantic anomaly detection catch a
kind of failure a confidence score could never flag in the first place, because every individual
decision can be scored with full confidence and still fail, when the scene being reasoned about is
not the scene the planner assumed it was.


% ---------- E10. STT section (sec:stt_lit): trim + lead with the solution + Fluister hedge ----------
% REPLACE THE WHOLE SECTION, from "\mysection{Speech recognition for local languages}..."
% down to (but not including) the following "\mysection{Research gap}...".

\mysection{Speech recognition for local languages}{Speech recognition for South African English and Afrikaans}
\label{sec:stt_lit}

Every method discussed so far assumes the instruction $d_i$ handed to the LLM planner is already
correct text. For a voice-commanded rover, that text comes from a speech-to-text stage first, and
an error there becomes an error in the planner's input before any clarification mechanism gets a
chance to act on it. South African English and Afrikaans are both under-represented in the
training data of general-purpose speech recognisers, which makes this an upstream risk specific to
a South African test environment.

\citet{kamper2012multiaccent} quantify this cost directly: across five South African English
accents, a single acoustic model shared across all accents performs about as well as one trained
per accent (around 81.5\% word accuracy either way), while a multi-accent model that shares
structure across accents but keeps accent-specific distinctions where they help reaches 82.8\%, a
statistically significant improvement over both. Black South African English was the one accent
that still benefited from being modelled separately, which the authors attribute to it being
acoustically the most distinct of the five.

\citet{jacobs2025afrikaans} report the same conclusion from a lower-resource direction:
fine-tuning Whisper on Afrikaans and isiXhosa child speech with only five minutes of transcribed
in-domain audio still gives a large accuracy gain, and that gain grows further when in-domain
adult speech and voice-conversion-based data augmentation are added.

Both results point the same way: off-the-shelf speech recognition degrades measurably on South
African English and Afrikaans, and the literature's fix is not more general-purpose training data
but a small amount of in-domain or in-accent fine-tuning, in the spirit of
\texttt{digiphyte/fluister-turbo}, a faster-whisper checkpoint fine-tuned specifically for South
African English and Afrikaans. Whether that advantage holds up against a larger general-purpose
model in this project's own setting is not guaranteed by either result above: a full-weight
multilingual Whisper model, run at a size the accent-specific checkpoint does not need to match,
can still win on raw transcription accuracy even without any accent-specific tuning. This is a
design constraint on the speech-to-text stage that sits upstream of every planning method
discussed in Sections~\ref{sec:llm_planners} to~\ref{sec:hallucination_lit}, not something those
methods can correct for after the fact, which is why this project evaluates the in-domain and
general-purpose options against each other directly rather than assuming the literature's
conclusion carries over unchanged.

